\chapter{Why a radiance field needs a resolution limit}
\label{ch:intro}

Suppose you are handed eighty photographs of an object and asked to produce a
picture of it from a viewpoint nobody stood at. This is the problem a radiance
field solves, and the last few years have made it look easy. It is worth
pausing, before any of the machinery, on what the photographs can and cannot
tell you.

They sample the scene, and they sample it unevenly. A surface close to the
cameras is covered by many more pixels than one far away. Some parts of the
object are seen from many directions, others only from a narrow range, and a few
may be seen from a single grazing angle or not at all. Different cameras may not
even have the same focal length, so their pixels do not correspond to the same
patch of surface. Whatever you reconstruct from such a set, there is a limit on
how fine the detail can be, and that limit belongs to the capture. It is not a
property of the representation you happen to store the answer in, and no choice
of representation can lift it.

The authors of NeRF understood this and acted on it. Their positional encoding
uses $L$ frequency bands, and they chose $L$ from, in their words, ``the maximum
frequency present in the sampled input images''. The instinct is right.
The instrument is as blunt as an instrument can be: one number, picked by hand,
applied to an entire scene.

3D Gaussian Splatting dropped even that. Its authors say plainly that no
regularisation is applied, and they mean it: a primitive may shrink without
bound, because nothing in the objective gives it a reason not to. As long as the
model is evaluated at the resolution it was fitted at, this does no visible harm, which is why it went unnoticed for so long. Change the
resolution and the damage appears at once. Section~\ref{sec:cost} puts a number
on it: $15.64$\,dB in the published figures, and \StmtArmBDrop\,dB in this
project's own reproduction.

Mip-Splatting restored a limit, and the restoration works; the numbers in
Chapter~\ref{ch:results} confirm it comfortably. But the limit it restores is a
single scalar per primitive, added in quadrature to all three axes of the
covariance. This thesis asks what that scalar ought to have been, and arrives at
the answer that it ought not to have been a scalar.

\section{The claim: the limit should be a matrix computed from the capture}

\keybox{thesis}{%
``The maximum frequency present in the sampled input images'' is not one number.
It falls with depth; it differs across the ray from along it; and it varies over
the sphere of directions. The accumulated sampling Fisher information is its
local, directional generalisation. Floored at the per-view Nyquist limit, it
reduces to Mip-Splatting exactly where the capture is well conditioned, and
departs from it precisely where the capture is not.}

\section{The same problem in signal-processing terms: reconstruction from non-uniform samples}

Stripped of its graphics vocabulary, the problem is reconstruction from
non-uniform samples without a matched prefilter --- a classical
signal-processing failure mode. The objects involved are familiar ones.

\begin{table}[htbp]
  \centering
  \caption{The same objects, in two vocabularies.}
  \small
  \begin{tabular}{p{0.42\linewidth}p{0.5\linewidth}}
    \toprule
    graphics term & signal-processing object \\
    \midrule
    Anti-aliasing filter (2D Mip filter) & Reconstruction prefilter matched to the display sampling rate \\
    3D smoothing filter & Prefilter matched to the \emph{acquisition} sampling rate \\
    Primitive covariance $\Sigma_k$ & Support of a basis function in a non-orthogonal expansion \\
    ``Too many small Gaussians'' & Over-parameterisation relative to the information in the measurements \\
    Floaters and depth artefacts & Fitting components along directions the data does not constrain \\
    Spherical-harmonic view dependence & Basis expansion whose Gram matrix may be rank-deficient \\
    \bottomrule
  \end{tabular}
\end{table}

The contribution is to compute the acquisition-side prefilter from the conditioning of the capture instead of choosing it by hand, and to keep it anisotropic
because the information is. \emph{An} anisotropic 3D filter is not itself new
--- AAA-Gaussians (2025) has one, derived from the view being rendered --- and
neither is a per-primitive second-order information matrix, which PUP 3D-GS
(2025) builds in order to prune. Section~\ref{sec:relatedwork} states the line of
separation precisely and Chapter~\ref{ch:discussion} returns to it; the short
form is that what is new here is the \emph{source} of the anisotropy and the
retention of its \emph{spectrum}.

\section{Contributions, and what each one rests on}

\begin{enumerate}[leftmargin=1.4em,itemsep=2pt]
  \item A derivation of the per-primitive, per-direction band-limit from the
        conditioning of the training capture --- an observability matrix
        accumulated from the same Jacobian the renderer already forms
        (\S\ref{sec:jacobian}--\S\ref{sec:accumulation}). The novelty is that
        the band-limit is a function of the cameras that \emph{saw} the
        primitive, not of the camera now rendering it, and that its whole spectrum is used
        (\S\ref{sec:relatedwork}).
  \item A proof that Mip-Splatting's scalar filter is the single-view,
        lateral-only, isotropised special case, with its constant $\sqrt{0.2}$
        identified as $s\,\sigma_{\mathrm{pix}} = 0.447$\,px
        (\S\ref{sec:proposition1}).
  \item Exact expressions for the anisotropy in three regimes: within one view,
        $\lambda_2/\lambda_1 = \cos^2\varphi$ off-axis (\S\ref{sec:offaxis});
        across two views, the triaxial spectrum
        $\{2, 2\cos^2(\theta/2), 2\sin^2(\theta/2)\}$
        (\S\ref{sec:anisotropy}); and over a distribution of views, the cap form of \S\ref{sec:capprediction}, the one a real protocol obeys and
        which the two-view form understates by $\sqrt{2}$. Together they predict
        where the method must help and where it must not.
  \item An exact-reduction theorem making parity on well-conditioned benchmarks
        structural (\S\ref{sec:proposition2}).
  \item An angular identifiability criterion for spherical-harmonic coefficients, which are $76\,\%$ of the model
  (Chapter~\ref{ch:method2}), together with the measurement that \textbf{refutes
  it}. At matched
        structure and matched budget it loses to magnitude pruning
        (\S\ref{sec:btwo}). The criterion still describes the geometry it was
        derived from; it is a poor rule for deciding what to keep, and is
        reported as such.
  \item A verified, executed reproduction of both baselines on free compute,
        with the twelve pre-flight assertions and the per-scale recovery that
        the shipped pipeline cannot produce (Chapters~\ref{ch:design} and
        \ref{ch:results}).
  \item A measured cost for the method, paired in one session on one card:
        $2.8\times$ peak VRAM and $1.32\times$ training time for a model whose
        primitive count matches the baseline's to $0.2\,\%$
        (\S\ref{sec:methodcost}). A reduction theorem says the output is
        identical; only a measurement says what the identity costs.
  \item Four defects in this project's own measurement, found by checks built
        for the purpose and reported with the corrections they forced: an arm
        that was not the baseline it was labelled (\S\ref{sec:g2}), a column
        averaged across two builds (\S\ref{sec:blend}), a reproduction error
        that was a seven-scene mean compared against an eight-scene published
        figure (\S\ref{sec:table1}), and an explanation for that error which
        was asserted for months and is now measured to be wrong
        (\S\ref{sec:densification}). Each was invisible in the numbers it
        produced; each is now covered by a check that fails the build.
\end{enumerate}
